\documentclass[aspectratio=169,11pt]{beamer}

% --- Theme ---
\usetheme{metropolis}
\usepackage{appendixnumberbeamer}

% --- Packages ---
\usepackage{fontspec}
\usepackage{minted}
\setminted{fontsize=\small,breaklines=true}
\usepackage{booktabs}
\usepackage{tikz}
\usepackage{fontawesome5}

% --- Colors ---
\definecolor{healthgreen}{RGB}{34,139,34}
\definecolor{alertred}{RGB}{220,53,69}
\setbeamercolor{alerted text}{fg=alertred}

% --- Metadata ---
\title{Module 8: Machine learning for clinical prediction}
\subtitle{Data analysis with Python for health specialists}
\author{Ya\'{e} Ulrich Gaba}
\institute{AIRINA Labs}
\date{2026}

\begin{document}

% ============================================================
\maketitle

% ============================================================
\begin{frame}{ML vs.\ traditional statistics}
\begin{columns}
\column{0.48\textwidth}
\textbf{Traditional statistics:}\\
``Which risk factors are associated with heart disease, and by how much?''

\vspace{4pt}
Goal = \alert{inference}\\
Tool = logistic regression\\
Output = odds ratios, p-values

\column{0.48\textwidth}
\textbf{Machine learning:}\\
``Given this patient's data, what is their probability of heart disease?''

\vspace{4pt}
Goal = \alert{prediction}\\
Tool = random forest, etc.\\
Output = predicted probability
\end{columns}

\vspace{8pt}
\begin{exampleblock}{In practice}
Many clinical studies use \textbf{both}: regression to understand the biology, ML to build the screening tool.
\end{exampleblock}
\end{frame}

% ============================================================
\begin{frame}{When NOT to use ML}
\begin{itemize}
  \item \textbf{Small datasets} ($< 200$ patients): ML overfits easily. Stick with logistic regression.
  \item \textbf{Regulatory approval:} FDA requires interpretable models for most clinical decision support.
  \item \textbf{Simple rules work:} If ``glucose $> 126$'' classifies 90\% of diabetics, you do not need a random forest.
\end{itemize}

\vspace{8pt}
\centering
\Large
\faIcon{balance-scale}\\[4pt]
{\normalsize Complexity must be \alert{justified by performance gain}.}
\end{frame}

% ============================================================
\begin{frame}[fragile]{Train/test split and cross-validation}
\begin{minted}{python}
from sklearn.model_selection import (train_test_split,
                                     cross_val_score)

# 70/30 split (stratified for class balance)
X_train, X_test, y_train, y_test = train_test_split(
    X, y, test_size=0.3, random_state=42, stratify=y)

# 5-fold cross-validation (more robust)
scores = cross_val_score(model, X, y, cv=5,
                         scoring="accuracy")
print(f"CV accuracy: {scores.mean():.3f} +/- {scores.std():.3f}")
\end{minted}

\begin{alertblock}{Golden rule}
\alert{Never} use test set results to choose your model. The test set is used \textbf{once}, at the very end.
\end{alertblock}
\end{frame}

% ============================================================
\begin{frame}[fragile]{Decision trees: clinical flowcharts}
\begin{minted}{python}
from sklearn.tree import DecisionTreeClassifier, plot_tree

dt = DecisionTreeClassifier(max_depth=4, random_state=42)
dt.fit(X_train, y_train)

# Visualize the tree
fig, ax = plt.subplots(figsize=(20, 10))
plot_tree(dt, feature_names=feature_cols,
          class_names=["No Disease", "Disease"],
          filled=True, rounded=True, ax=ax)
\end{minted}

\begin{exampleblock}{Clinical relevance}
Decision trees produce \textbf{clinical decision rules} --- flowcharts a physician can follow at the bedside. The HEART score, Wells criteria, and Ottawa ankle rules are conceptually similar.
\end{exampleblock}
\end{frame}

% ============================================================
\begin{frame}[fragile]{Random forests: ensemble power}
\begin{minted}{python}
from sklearn.ensemble import RandomForestClassifier

rf = RandomForestClassifier(n_estimators=200, max_depth=6,
                            random_state=42)
rf.fit(X_train, y_train)
y_pred = rf.predict(X_test)

print(classification_report(y_test, y_pred,
      target_names=["No Disease", "Disease"]))
\end{minted}

\vspace{4pt}
A random forest = \textbf{200 decision trees}, each on a random subset. Final prediction = majority vote. More accurate, less interpretable.
\end{frame}

% ============================================================
\begin{frame}[fragile]{ROC curves and AUC}
\begin{minted}{python}
from sklearn.metrics import roc_curve, roc_auc_score

y_prob = rf.predict_proba(X_test)[:, 1]
fpr, tpr, _ = roc_curve(y_test, y_prob)
auc = roc_auc_score(y_test, y_prob)

ax.plot(fpr, tpr, label=f"RF (AUC = {auc:.3f})")
ax.plot([0, 1], [0, 1], "k--", label="Random (0.500)")
\end{minted}

\begin{table}
\begin{tabular}{ll}
\toprule
\textbf{AUC} & \textbf{Discrimination} \\
\midrule
0.5 & No better than a coin flip \\
0.7--0.8 & Acceptable \\
0.8--0.9 & Excellent \\
$>$0.9 & Outstanding (rare in practice) \\
\bottomrule
\end{tabular}
\end{table}
\end{frame}

% ============================================================
\begin{frame}[fragile]{Feature importance}
\textbf{Q:} Which risk factors matter most for prediction?

\begin{minted}{python}
importances = rf.feature_importances_
imp_df = pd.DataFrame({
    "Feature": feature_cols,
    "Importance": importances
}).sort_values("Importance", ascending=True)

ax.barh(imp_df["Feature"], imp_df["Importance"])
ax.set_xlabel("Feature importance (Gini)")
\end{minted}

\begin{alertblock}{Correlation $\neq$ Causation}
Feature importance tells you what is useful for \alert{prediction}, not what \alert{causes} the outcome. A proxy variable can rank highly.
\end{alertblock}
\end{frame}

% ============================================================
\begin{frame}[fragile]{Calibration: do probabilities match reality?}
\begin{minted}{python}
from sklearn.calibration import calibration_curve
from sklearn.metrics import brier_score_loss

prob_true, prob_pred = calibration_curve(y_test, y_prob,
                                         n_bins=8)
ax.plot(prob_pred, prob_true, "o-", label="Random Forest")
ax.plot([0, 1], [0, 1], "k--", label="Perfect")

brier = brier_score_loss(y_test, y_prob)
print(f"Brier score: {brier:.4f}")  # lower = better
\end{minted}

\vspace{4pt}
If a model says ``70\% risk,'' do 70\% of those patients actually have the disease? \alert{Calibration is critical for clinical deployment.}
\end{frame}

% ============================================================
\begin{frame}[fragile]{Overfitting: the central danger}
\begin{minted}{python}
# Unlimited depth -> memorizes training data
dt_overfit = DecisionTreeClassifier()  # no max_depth
dt_overfit.fit(X_train, y_train)

print(f"Train acc (overfit): {dt_overfit.score(X_train, y_train):.3f}")
print(f"Test acc (overfit):  {dt_overfit.score(X_test, y_test):.3f}")
# 1.000 vs 0.72 -> overfitting!
\end{minted}

\vspace{4pt}
\textbf{Signs:} training $\gg$ test accuracy, wild CV fold variation, more features hurt.\\
\textbf{Cure:} more data, simpler models, regularization, cross-validation.
\end{frame}

% ============================================================
\begin{frame}{Ethics: fairness in clinical ML}
Before deploying any clinical ML model, ask:

\begin{itemize}
  \item \textbf{Who is in the training data?} 90\% white males $\to$ may fail for women and minorities.
  \item \textbf{Does performance differ across subgroups?} Check AUC by sex, race, age.
  \item \textbf{What are the consequences of errors?} A missed cardiac event in a young woman is dangerous.
  \item \textbf{Is the model transparent?} Clinicians and patients deserve explanations.
\end{itemize}

\begin{alertblock}{Real-world example}
In 2019, a widely-used algorithm was found to be biased against Black patients --- it used healthcare \emph{costs} as a proxy for \emph{needs}, systematically underestimating care requirements.
\end{alertblock}
\end{frame}

% ============================================================
\begin{frame}{What you can do after this module}
\begin{enumerate}
  \item Split data properly and use cross-validation
  \item Train decision trees, random forests, and logistic regression
  \item Compare models using ROC curves and AUC
  \item Identify top predictive features from a random forest
  \item Assess calibration and detect overfitting
  \item Audit models for fairness across demographic subgroups
\end{enumerate}

\vspace{12pt}
\centering
\Large \alert{Next: Module 9 --- Geospatial and temporal health data}
\end{frame}

% ============================================================
\begin{frame}{Exercises (Hour 3)}
\begin{enumerate}
  \item \textbf{Model comparison:} Train logistic regression, decision tree, and random forest; compare CV AUC
  \item \textbf{Decision tree:} Visualize a depth-3 tree; can a clinician follow it?
  \item \textbf{Top features:} Train RF with only top 5 features; compare AUC to full model
  \item \textbf{Mini-project:} Build a complete heart disease prediction pipeline with ROC curves, feature importance, calibration plot, and fairness audit
\end{enumerate}
\end{frame}

\end{document}
